\documentclass[11pt]{article}
\usepackage[a4paper,margin=26mm]{geometry}
\usepackage{amsmath,amssymb,mathtools}
\usepackage{unicode-math}
\setmainfont{Latin Modern Roman}
\setmathfont{Latin Modern Math}
\usepackage{xurl,hyperref}
\hypersetup{hidelinks,pdftitle={Problem statement}}
\newcommand{\sref}[2]{\hyperlink{source:#1:#2}{[S#2]}}
\newcommand{\cataloguescope}[1]{#1}
\begin{document}
\section{Leaf spaces and transverse structures — Question 8.10}\label{statement}

\subsection{Definitions and mathematical statement}

\paragraph{Shared notation.}$\mathbb N=\{1,2,\ldots\}$, and $\mathbb Z,\mathbb Q,\mathbb R,\mathbb C$ denote the integers, rationals, reals and complex numbers. Any other conventions are specified locally.


A codimension-one foliation $\mathcal F$ decomposes a $3$-manifold $M$ into surface leaves, locally as the horizontal slices of a product chart. In the terminology of the source, it is taut if a closed curve transverse to $\mathcal F$ meets every leaf. A minimal foliation means that every leaf is dense; a minimal surface means a surface with vanishing mean curvature for the specified Riemannian metric. These two uses of minimal are distinct. A manifold is atoroidal if it has no incompressible embedded torus other than a boundary-parallel one. Write $\widetilde M$ for its universal cover, and put a tilde on a foliation or lamination for its lift. For a taut foliation with hyperbolic lifted leaves, a universal circle is a circle $S^1_{\mathrm{univ}}$ with a $\pi_1(M)$-action and continuously varying monotone maps $\phi_\lambda:S^1_{\mathrm{univ}}\to S^1_\infty(\lambda)$ to leaf ideal circles. Monotone means point fibers are contractible. The maps are equivariant: $\phi_{g\lambda}\circ\rho(g)=g\circ\phi_\lambda$. The core for a leaf is the complement of the open intervals collapsed by its map. For incomparable lifted leaves, each core lies in one gap of the other leaf. A minimal circle action has every orbit dense.
\paragraph{Statement recorded in the supplied source.}
Determine the optimal analytic regularity obtainable for the $\pi_1(M)$-action on $S^1_{\mathrm{univ}}$, allowing a change of circle coordinates. In particular compare Lipschitz, differentiable and more regular realizations of these actions. The target concerns the universal-circle action together with its foliation origin, rather than arbitrary circle actions of the group. \sref{statement}{2}
\subsection{Short English statement}
Determine how regular an action on a foliation universal circle can be made by changing circle coordinates.
\subsection{Sources}
\begin{description}
\item[\hypertarget{source:statement:1}{S1}] ULAM.AI and the UnsolvedMath Contributors, \emph{UnsolvedMath: A Curated Collection of Open Mathematics Problems}, catalogue entry. CC BY 4.0. \url{https://huggingface.co/datasets/ulamai/UnsolvedMath}.
\item[\hypertarget{source:statement:2}{S2}] Danny Calegari, \emph{Problems in foliations and laminations of 3-manifolds}, version 0.78 (2002), Question 8.10; Definitions 1.1--1.10 and local definitions and remarks. \url{https://arxiv.org/abs/math/0209081v1}.
\end{description}
\label{end:statement}
\end{document}
